% 24. ТВ 3.pdf, с. 2–3 (формулировка и доказательство ЗБЧ, обе оценки);
% с. 6 (обозначение сходимости по вероятности для выборочного среднего).
\question{Неравенства Маркова и Чебышева. Закон больших чисел}

\begin{theorem}[\topicmark{Неравенство Маркова}Неравенство Маркова]
Если $\xi\geq0$ почти наверное и $\EE\xi<\infty$, то при любом $\eps>0$
\[
 \Pr\{\xi\geq\eps\}\leq\frac{\EE\xi}{\eps}.
\]
\end{theorem}
\begin{Proof}
Положим $A=\{\xi\geq\eps\}$. Разложим ожидание по событиям $A$ и $\overline A$.
Оба слагаемых неотрицательны, а на $A$ величина $\xi$ не меньше $\eps$. Поэтому
\[
 \EE\xi=\EE(\xi I_A)+\EE(\xi I_{\overline A})
 \geq\EE(\xi I_A)\geq\eps\,\EE I_A=\eps\,\Pr(A).
\]
Делим на $\eps>0$.
\end{Proof}

\begin{theorem}[\topicmark{Неравенство Чебышева}Неравенство Чебышева]
Если у $\xi$ конечная дисперсия, то при любом $\eps>0$
\[
 \Pr\{|\xi-\EE\xi|\geq\eps\}\leq\frac{\DD\xi}{\eps^2}.
\]
\end{theorem}
\begin{Proof}
Применим неравенство Маркова к $(\xi-\EE\xi)^2$ с порогом $\eps^2$:
\[
 \Pr\{|\xi-\EE\xi|\geq\eps\}
 =\Pr\{(\xi-\EE\xi)^2\geq\eps^2\}
 \leq\frac{\EE(\xi-\EE\xi)^2}{\eps^2}=\frac{\DD\xi}{\eps^2}.
\]
\end{Proof}

\begin{definition}[\topicmark{Сходимость по вероятности}Сходимость по вероятности\editorialmark]
Пусть $\xi_1,\xi_2,\ldots$ и $\xi$ заданы на одном вероятностном
пространстве. Последовательность $\xi_n$ \textbf{сходится по вероятности}
к $\xi$, если для каждого $\eps>0$
\[
    \lim_{n\to\infty}\Pr\{|\xi_n-\xi|>\eps\}=0.
\]
Обозначение: $\xi_n\xrightarrow{\Pr}\xi$. При постоянном пределе $\xi=a$
под знаком вероятности получаем событие $\{|\xi_n-a|>\eps\}$.
\end{definition}

\Needspace{16\baselineskip}
\begin{theorem}[\topicmark{Закон больших чисел}Закон больших чисел]
Пусть $\xi_1,\xi_2,\ldots$ независимы, одинаково распределены,
$\EE\xi_1=a$ и $\DD\xi_1=\sigma^2<\infty$. Тогда их среднее
$\overline\xi_n=(\xi_1+\cdots+\xi_n)/n$ сходится по вероятности к $a$:
\[
 \overline\xi_n\xrightarrow{\Pr}a,\qquad
 \Pr\{|\overline\xi_n-a|>\eps\}\longrightarrow 0
 \quad\text{при каждом }\eps>0.
\]
\end{theorem}
\begin{Proof}
По линейности $\EE\overline\xi_n=a$. Независимость слагаемых даёт
\[
 \DD\overline\xi_n
 =\frac{\DD(\xi_1+\cdots+\xi_n)}{n^2}
 =\frac{n\sigma^2}{n^2}=\frac{\sigma^2}{n}.
\]
По неравенству Чебышева
\[
 \Pr\{|\overline\xi_n-a|>\eps\}\leq
 \Pr\{|\overline\xi_n-a|\geq\eps\}
 \leq\frac{\sigma^2}{n\eps^2}\xrightarrow[n\to\infty]{}0.
\]
\end{Proof}

% 25. ТВ 3.pdf, с. 2, 6. В программе требуется формулировка ЦПТ.
\question{Сходимость по распределению и центральная предельная теорема}

\begin{definition}[\topicmark{Сходимость по распределению}Сходимость по распределению]
Пишут $\eta_n\xrightarrow{d}\eta$, если для каждой точки $x$, в которой
$F_\eta$ непрерывна, выполнено
\[
 F_{\eta_n}(x)\xrightarrow[n\to\infty]{}F_\eta(x).
\]
\end{definition}

% ТВ 3.pdf, с. 6: пример, в котором предел имеет скачок в нуле.
\begin{example}[Что происходит в точке разрыва]
Пусть $\eta_n$ принимает значения $0$ и $1/n$ с вероятностями $1/2$,
а $\eta=0$ почти наверное.

\smallskip
При $x<0$ обе функции распределения равны нулю.
При каждом $x>0$ начиная с некоторого $n$ обе равны единице.

\smallskip
В точке $x=0$ имеем $F_{\eta_n}(0)=1/2$ и $F_\eta(0)=1$.
Это точка разрыва $F_\eta$, поэтому условие сходимости по распределению выполнено.
\end{example}

\begin{theorem}[\topicmark{Центральная предельная теорема}Центральная предельная теорема]
Пусть $\xi_1,\xi_2,\ldots$ независимы и одинаково распределены,
$\EE\xi_1=a$, $0<\DD\xi_1=\sigma^2<\infty$. Если $Z\sim N(0;1)$, то
\[
 \frac{\xi_1+\cdots+\xi_n-na}{\sigma\sqrt n}\xrightarrow{d}Z.
\]
\end{theorem}

% 26. ТВ 3.pdf, с. 6–8; Тервер(Лекция)69.pdf, с. 4; lecture_17_expectation.tex:108–113.
\question{Характеристическая функция и её свойства}
\begin{definition}[\topicmark{Определение}Характеристическая функция]
Характеристической функцией случайной величины $\xi$ называется функция
\[
 \varphi_\xi(t)=\EE e^{it\xi},\qquad t\in\R,\quad i^2=-1.
\]
\end{definition}

\begin{theorem}[\topicmark{Существование}Существование характеристической функции]
У любой случайной величины $\xi$ характеристическая функция существует при всех $t\in\R$.
\end{theorem}
\begin{Proof}
По формуле Эйлера
\[
 \varphi_\xi(t)=\EE\cos(t\xi)+i\,\EE\sin(t\xi).
\]
Поскольку $|\cos(t\xi)|\leq1$ и $|\sin(t\xi)|\leq1$, оба математических ожидания существуют.
\end{Proof}

\begin{theorem}[\topicmark{Единственность}Теорема единственности]
Величины $\xi$ и $\eta$ имеют одинаковое распределение
$\Leftrightarrow$ $\varphi_\xi(t)=\varphi_\eta(t)$ при всех $t\in\R$.
\end{theorem}

\begin{theorem}[\topicmark{Сумма независимых величин}Сумма независимых величин]
Если $\xi$ и $\eta$ независимы, то
$\varphi_{\xi+\eta}(t)=\varphi_\xi(t)\varphi_\eta(t)$ для всех $t\in\R$.
\end{theorem}
\begin{Proof}
Из независимости $\xi$ и $\eta$ следует независимость $e^{it\xi}$ и $e^{it\eta}$. Поэтому
\[
 \varphi_{\xi+\eta}(t)=\EE e^{it(\xi+\eta)}
 =\EE(e^{it\xi}e^{it\eta})
 =\EE e^{it\xi}\,\EE e^{it\eta}=\varphi_\xi(t)\varphi_\eta(t).
\]
\end{Proof}

% ТВ 3.pdf, с. 6–7.
\begin{example}[Сумма независимых пуассоновских величин]
Для $\xi\sim\operatorname{Poiss}(\lambda)$ по определению ожидания
\[
 \varphi_\xi(t)=e^{-\lambda}\sum_{k=0}^\infty\frac{(\lambda e^{it})^k}{k!}
 =\exp\{\lambda(e^{it}-1)\}.
\]
У независимых пуассоновских величин с параметрами $\lambda_1,\lambda_2$
произведение характеристических функций равно $\exp\{(\lambda_1+\lambda_2)(e^{it}-1)\}$.
По теореме единственности их сумма имеет распределение $\operatorname{Poiss}(\lambda_1+\lambda_2)$.
\end{example}

\Needspace{10\baselineskip}
\begin{theorem}[\topicmark{Линейное преобразование}Линейное преобразование]
Пусть $\eta=a\xi+b$, где $a,b\in\R$. Тогда
$\varphi_\eta(t)=e^{itb}\varphi_\xi(at)$.
\end{theorem}
\begin{Proof}
Множитель $e^{itb}$ постоянен. По определению
\[
 \varphi_\eta(t)=\EE e^{it(a\xi+b)}
 =e^{itb}\EE e^{i(at)\xi}=e^{itb}\varphi_\xi(at).
\]
\end{Proof}

% ТВ 3.pdf, с. 7: вычисление для N(0;1), затем применение линейного преобразования.
\begin{example}[Характеристическая функция нормального распределения]
Пусть $Z\sim N(0;1)$. В интеграле для $\varphi_Z(t)$ слагаемое с $\sin(tx)$
обращается в нуль по нечётности.

\smallskip
Дифференцируя оставшийся интеграл и
интегрируя по частям, получаем
\[
 \varphi_Z'(t)=-\frac1{\sqrt{2\pi}}\int_{\R} x\sin(tx)e^{-x^2/2}\,\diff x
 =-\frac{t}{\sqrt{2\pi}}\int_{\R}\cos(tx)e^{-x^2/2}\,\diff x
 =-t\varphi_Z(t).
\]
Из $\varphi_Z(0)=1$ следует $\varphi_Z(t)=e^{-t^2/2}$.
Поэтому для $\eta=\sigma Z+a$ имеем
$\varphi_\eta(t)=e^{ita}\varphi_Z(\sigma t)=e^{ita-\sigma^2t^2/2}$.
\end{example}

% ТВ 3.pdf, с. 8: версия с заданной предельной случайной величиной.
\begin{theorem}[\topicmark{Теорема Леви}Теорема непрерывности Леви]
Для случайных величин $\xi_n$ и $\xi$ справедлива эквивалентность
\[
 \xi_n\xrightarrow{d}\xi
 \quad\Longleftrightarrow\quad
 \varphi_{\xi_n}(t)\xrightarrow[n\to\infty]{}\varphi_\xi(t)
 \quad\text{для каждого }t\in\R.
\]
\end{theorem}
